\section{Portfolio Evaluation}
\label{sec:portfolio_implications}

The portfolio analysis applies the Fragility Score using information available at each quarterly formation date. The primary equal-weighted comparison contrasts all eligible stocks with a portfolio excluding the most fragile decile. Decile~1 and Decile~10 serve as diagnostic portfolios showing the score's economic separation.

\subsection{Primary Long-Only Portfolios}
\label{sec:primary_portfolio_results}

Table~\ref{tab:portfolio-event-windows} summarizes the eight quarterly 63-day holding windows. Excluding Decile~10 reduces average maximum drawdown, downside volatility, and the worst five-day loss relative to the full equal-weighted universe. The contrast between Decile~1 and Decile~10 is substantially larger, confirming that the score separates stocks with very different subsequent downside outcomes. The corresponding graphical comparison is reported in Appendix~\ref{app:portfolio_window_risk}.

\begin{table}[!htbp]
\centering
\footnotesize
\caption{Mean performance across quarterly 63-day portfolio windows.}
\label{tab:portfolio-event-windows}
\begin{tabular}{lrrrr}
\toprule
Portfolio & Max. drawdown & Downside vol. & Worst 5-day loss & Cum. return \\
\midrule
All eligible stocks & 9.6\% & 13.5\% & 6.3\% & 3.2\% \\
Exclude Decile 10 & 8.6\% & 12.7\% & 5.7\% & 4.5\% \\
Decile 1 & 4.6\% & 6.9\% & 3.2\% & 2.2\% \\
Decile 10 & 24.4\% & 27.1\% & 13.6\% & -8.9\% \\
\bottomrule
\end{tabular}
\end{table}


\noindent
The continuous-strategy results in Table~\ref{tab:portfolio-performance} and Figure~\ref{fig:portfolio-strategy-wealth} point in the same direction. The screened portfolio improves on the matched equal-weighted universe in return, volatility, Sharpe ratio, and maximum drawdown. It nevertheless underperforms the external CRSP total-market benchmark, which has both higher returns and a better downside profile over the same period.

\begin{table}[!htbp]
\centering
\footnotesize
\caption{Continuous Fragility Score portfolio performance.}
\label{tab:portfolio-performance}
\begin{tabular}{lrrrrr}
\toprule
Strategy & Cum. return & Ann. return & Ann. vol. & Sharpe & Max. drawdown \\
\midrule
All eligible stocks & 16.2\% & 7.8\% & 20.4\% & 0.24 & 26.5\% \\
Exclude Decile 10 & 29.5\% & 13.8\% & 19.5\% & 0.52 & 24.1\% \\
CRSP total market & 41.4\% & 19.0\% & 16.4\% & 0.85 & 19.1\% \\
\bottomrule
\end{tabular}
\end{table}


\begin{figure}[!htbp]
    \centering
    \includegraphics[width=0.82\textwidth]{04_11_fragility_strategy_wealth.pdf}
    \caption[Cumulative wealth of the primary portfolio strategies]{Cumulative wealth of the full equal-weighted universe, the portfolio excluding Decile~10, and the CRSP total-market index.}
    \label{fig:portfolio-strategy-wealth}
\end{figure}

\noindent
Figure~\ref{fig:portfolio-strategy-alpha} presents the factor-adjusted evidence. The screened portfolio has approximately zero FF5-plus-momentum alpha, whereas the screened-minus-unscreened portfolio has a positive alpha whose 95\% confidence interval excludes zero. This indicates a benefit relative to the matched internal universe, not superiority to the broad market. Complete regression statistics appear in Appendix~\ref{app:primary_portfolio_alpha}.

\begin{figure}[!htbp]
    \centering
    \includegraphics[width=0.76\textwidth]{04_11_fragility_strategy_alpha.pdf}
    \caption[Factor-adjusted alpha of the primary portfolios]{Annualized FF5-plus-momentum alpha with 95\% confidence intervals for the full, screened, and screened-minus-full portfolios.}
    \label{fig:portfolio-strategy-alpha}
\end{figure}

\subsection{Long-Only Portfolio Robustness}
\label{sec:portfolio_robustness_results}

The long-only robustness analysis tests sensitivity to the exclusion threshold and weighting rule. Table~\ref{tab:portfolio-robustness} shows that excluding the top 20\% produces results similar to excluding Decile~10 under equal weighting. Under value weighting, however, excluding Decile~10 has almost no effect. The score's benefit is therefore concentrated among smaller stocks and depends on the weighting rule, consistent with the firm-size error pattern in Section~\ref{sec:selected_model_errors}. A graphical comparison appears in Appendix~\ref{app:portfolio_robustness_figure}.

\begin{table}[!htbp]
\centering
\footnotesize
\caption{Long-only portfolio robustness results.}
\label{tab:portfolio-robustness}
\begin{tabular}{lrrrr}
\toprule
Strategy & Ann. return & Ann. vol. & Sharpe & Max. drawdown \\
\midrule
Equal weight: all stocks & 7.8\% & 20.4\% & 0.24 & 26.5\% \\
Equal weight: exclude Decile 10 & 13.8\% & 19.5\% & 0.52 & 24.1\% \\
Equal weight: exclude top 20\% & 13.4\% & 18.5\% & 0.51 & 23.4\% \\
Value weight: all stocks & 19.3\% & 16.6\% & 0.86 & 19.4\% \\
Value weight: exclude Decile 10 & 19.2\% & 16.6\% & 0.86 & 19.4\% \\
\bottomrule
\end{tabular}
\end{table}


\FloatBarrier

\subsection{Long--Short Extensions}
\label{sec:long_short_results}

The long--short analysis examines the return spread between the least- and most-fragile stocks. DFRAG-10 uses the extreme deciles, DFRAG-20 uses the two lowest and highest deciles, and volatility-managed DFRAG-10 scales exposure using lagged realized volatility. Table~\ref{tab:long-short-performance} and Figure~\ref{fig:long-short-wealth} compare the three specifications. DFRAG-10 produces the highest return but is volatile and concentrated, whereas volatility management reduces volatility and maximum drawdown while retaining positive returns.

\begin{table}[!htbp]
\centering
\footnotesize
\caption{Exploratory long--short Fragility Score strategies.}
\label{tab:long-short-performance}
\begin{tabular}{lrrrr}
\toprule
Strategy & Ann. return & Ann. volatility & Sharpe & Max. drawdown \\
\midrule
DFRAG-10 & 55.7\% & 34.6\% & 1.45 & 41.4\% \\
DFRAG-20 & 16.0\% & 29.6\% & 0.65 & 42.5\% \\
DFRAG-10: volatility managed & 28.2\% & 17.3\% & 1.52 & 21.1\% \\
\bottomrule
\end{tabular}
\end{table}


\begin{figure}[H]
    \centering
    \includegraphics[width=0.76\textwidth]{04_12_long_short_fragility_wealth.pdf}
    \caption[Cumulative wealth of the exploratory long--short strategies]{Cumulative wealth of DFRAG-10, DFRAG-20, and volatility-managed DFRAG-10. The horizontal line at one denotes zero return.}
    \label{fig:long-short-wealth}
\end{figure}

\noindent
The FF5-plus-momentum alpha estimate for DFRAG-10 is positive and its 95\% confidence interval excludes zero, whereas the DFRAG-20 and volatility-managed estimates are less precise. These estimates are shown in Figure~\ref{fig:long-short-alpha}. They should nevertheless be interpreted cautiously because the evaluation period is short and the strategies require substantial short-side implementation.

\begin{figure}[!htbp]
    \centering
    \includegraphics[width=0.76\textwidth]{04_12_long_short_fragility_alpha.pdf}
    \caption[Factor-adjusted alpha of the exploratory long--short strategies]{Annualized FF5-plus-momentum alpha with 95\% confidence intervals for DFRAG-10, DFRAG-20, and volatility-managed DFRAG-10.}
    \label{fig:long-short-alpha}
\end{figure}

\noindent
Table~\ref{tab:long-short-implementation} makes this limitation explicit. Average one-way turnover is high, particularly for volatility-managed DFRAG-10. Applying a 50-basis-point one-way trading cost reduces performance but does not eliminate the gross return in this short sample. The corresponding results under 10-, 25-, and 50-basis-point assumptions are reported in Appendix~\ref{app:long_short_cost_sensitivity}. These calculations still omit several material frictions, including borrowing availability, stock-loan fees, and market impact. The long--short evidence is therefore exploratory and does not support a tradable-factor interpretation.

\begin{table}[!htbp]
\centering
\footnotesize
\caption{Mean one-way turnover and annualized gross and net returns of the long--short strategies.}
\label{tab:long-short-implementation}
\begin{tabular}{lrrr}
\toprule
Strategy & Mean turnover & Gross return & Net return (50 bps) \\
\midrule
DFRAG-10 & 96.2\% & 55.7\% & 52.7\% \\
DFRAG-20 & 83.3\% & 16.0\% & 14.0\% \\
DFRAG-10: volatility managed & 130.7\% & 28.2\% & 24.9\% \\
\bottomrule
\end{tabular}
\end{table}


\noindent
Overall, the Fragility Score has clear economic ordering power and modestly improves the downside profile of the matched equal-weighted portfolio. Its portfolio value is qualified by weighting sensitivity, implementation costs, and the stronger performance of the external market benchmark.
